\subsection*{الف}
با توجه به اینکه تصاویر حقیقی هستند:
\begin{itemize}
\item اگر $f(x,y)$ نسبت به مبدا فرد باشد، $F(u,v)$ کاملا موهومی است.
\item اگر $f(x,y)$ نسبت به مبدا زوج باشد، $F(u,v)$ کاملا حقیقی است.
\item $F(0,0)$ برابر مجموع مقادیر تصویر است.
\end{itemize}
نتیجه‌ها:
\begin{itemize}
\item (i) شکل (e)
\item (ii) شکل‌های (a)، (b)، (c)، (f)، (h)
\item (iii) هیچ‌کدام
\item (iv) شکل‌های (e)، (h)
\item (v) شکل (f)
\end{itemize}

\subsection*{ب}
برای هر تصویر، تبدیل فوریه گسسته به صورت زیر است:
$$
F(u,v)=\sum_{x=0}^{M-1}\sum_{y=0}^{N-1} f(x,y)\,e^{-j2\pi\left(\frac{ux}{M}+\frac{vy}{N}\right)}
$$
و قدرمطلق و فاز به شکل زیر محاسبه می‌شوند:
$$
|F|=\sqrt{(\Re F)^2+(\Im F)^2}, \qquad
\angle F=\arctan\left(\frac{\Im F}{\Re F}\right)
$$
برای نمایش بهتر طیف، از انتقال مرکز با \lr{fftshift} و نمایش لگاریتمی \lr{log(1+|F|)} استفاده می‌کنیم. روابط کلی بین دامنه مکانی و فرکانسی:
\begin{itemize}
\item ساختارهای تکراری، قله‌های تیز در فرکانس ایجاد می‌کنند.
\item لبه‌های جهت‌دار، انرژی را روی خط عمود بر جهت لبه متمرکز می‌کنند.
\item نواحی صاف و بزرگ، انرژی کم‌فرکانس نزدیک مرکز طیف دارند.
\item بافت‌های نامنظم، طیف پخش و گسترده می‌سازند.
\item ساختارهای مدور، تقارن دایره‌ای در طیف ایجاد می‌کنند.
\end{itemize}